% Generated draft: requires mathematical and source review.
\documentclass[11pt]{article}
\usepackage[margin=25mm]{geometry}
\usepackage{fontspec}
\setmainfont{Latin Modern Roman}
\usepackage{amsmath,amssymb,mathtools,unicode-math}
\setmathfont{Latin Modern Math}
\usepackage{enumitem,needspace,xurl,hyperref}
\setlength{\parindent}{0pt}
\setlength{\parskip}{5pt}
\setlength{\emergencystretch}{3em}
\newcommand{\N}{\mathbb{N}}
\newcommand{\Z}{\mathbb{Z}}
\newcommand{\Q}{\mathbb{Q}}
\newcommand{\R}{\mathbb{R}}
\newcommand{\C}{\mathbb{C}}
\newcommand{\F}{\mathbb{F}}
\newcommand{\A}{\mathbb{A}}
\newcommand{\PP}{\mathbb{P}}
\newcommand{\E}{\mathbb{E}}
\newcommand{\eps}{\varepsilon}
\newcommand{\dd}{\,\mathrm{d}}
\begin{document}
\section*{UnsolvedMath ID: 11000068}
\textbf{Source problem code:} catalogue entry
\subsection{Definitions and mathematical statement}
Let \(\Sigma_g\) be a closed oriented genus-\(g\) surface. For \(L\geq2\), the full level subgroup is
\[\operatorname{Mod}_g[L]=\ker\bigl(\operatorname{Mod}_g\to\operatorname{Sp}(H_1(\Sigma_g;\Z/L\Z))\bigr).\]
The Picard-number-one conjecture for level structures is printed in the extraction as \(H2(\operatorname{Mod}_g[L];\Q)=\Q\) for \(g\geq3\), followed by the task of computing the integral \(H2\) for every \(g\geq3,L\geq2\). The index position is lost in the supplied text. Under the homological reading, the precise targets are
\[\forall g\geq3\ \forall L\geq2,\qquad H_2(\operatorname{Mod}_g[L];\Q)\cong\Q,\]
and determination, as abelian groups, of
\[H_2(\operatorname{Mod}_g[L];\Z)\quad(g\geq3,L\geq2).\]
If the source instead prints superscripts, these are cohomology targets \(H^2\). Rational rank-one statements correspond through the universal coefficient theorem, but the integral homology and cohomology requests cannot be interchanged: integral cohomology also reflects \(\operatorname{Ext}^1(H_1,\Z)\). Thus this typographic point requires confirmation, not a silent choice based only on the Picard name.

The source reports Harer [Har2] for level \(L=1\), and says the Picard-group generalization appears as Question 7.12 in [HL]. It reports a claimed \(L=2\) proof in [Fo] with an apparent error, and says even \((g,L)=(3,2)\) was open at that time. These are historical qualifications, not a present assertion that the same case remains open.

It proposes an approach for \(g\geq4\): \(\operatorname{Mod}_g[L]\) is finite index in the finitely presented \(\operatorname{Mod}_g\), hence finitely presented, and the finite symplectic group \(\operatorname{Sp}(2g,\Z/L\Z)\) has explicit structure. A subsequent question is announced but not provided, so no particular presentation or computation is inserted. The level group is a mapping-class subgroup; the summary's analogy with locally symmetric/Shimura Picard conjectures does not define this group or verify its homology.
\subsection{Short English statement}
Prove rational rank one in degree two for level mapping class groups and compute the integral degree-two groups in every stated genus and level. The extraction loses homology-versus-cohomology notation, and the historical level-two claim/error must remain qualified.
\subsection{Sources}
\begin{itemize}
\item Local UnsolvedMath registry, permanent ID \(11000068\), exact source code \texttt{catalogue entry}.
\item Farb, \emph{Problems on Mapping Class Groups and Related Topics}, Conjecture 5.24, PDF page 52: \url{https://www.math.uchicago.edu/~farb/papers/mcgbook.pdf}. Registry access 2026-07-29; pdf-font-structure extraction.
\item Harer [Har2] for \(L=1\), [HL] Question 7.12 for the Picard-group question, and [Fo] for the claimed level-two proof with an apparent error. Their full bibliographic entries are not supplied.
\end{itemize}
\end{document}
